\section{Extended Data figures and the claims ledger}
\label{ann:display}

\subsection{Extended Data figures}

\begin{figure}[hbtp]
  \centering
  \includegraphics[width=0.95\linewidth]{fig1_overview.png}
  \caption{\textbf{Study design, invariants and the four questions.} The harmonized
  three-cohort FACE baseline (bipolar disorder, schizophrenia, depression;
  $N=9{,}013$) is reorganized around diagnosis-agnostic axes by one global,
  missingness-aware Bayesian bifactor / exploratory structural-equation model fit to
  each patient's observed cells only, under three core modelling principles (no
  imputation; diagnosis as metadata; baseline defines the map and follow-up validates
  it). The fitted map is then interrogated through four dependent, increasingly
  demanding questions---exists, organizes, persists, predicts---with
  positive and null results reported alike.}
  \label{fig:design}
\end{figure}

\begin{figure}[p]
  \centering
  \includegraphics[height=0.86\textheight]{edfig_full_atlas.png}
  \caption{\textbf{The complete posterior loading atlas.} As Figure~\ref{fig:atlas}a but
  showing \emph{every} modelled indicator (all $109$ items: $88$ continuous and $21$ explicit) grouped
  by home factor, with no per-block cap. Dot size and colour encode the absolute posterior
  median loading; outlined dots have a $95\%$ credible interval excluding zero (heavier ring
  $=$ home-factor anchor). Three data-supported cross-loadings render as outlined dots in the
  cognition column---childhood adversity (CTQ-37, $-0.094$) and two poor-sleep items
  (PSQI latency $+0.057$, PSQI daytime dysfunction $-0.070$), each with a $95\%$ credible
  interval excluding zero. The diagonal block structure---each clinical block loading
  cleanly on its own specific axis, with a sparse general-burden ($\Gfac$) column---is the
  full anatomy of the measurement model summarised in the main-text Figure~\ref{fig:atlas}.}
  \label{fig:fullatlas}
\end{figure}

\begin{figure}[hbtp]
  \centering
  \includegraphics[width=0.72\linewidth]{edfig_consort.png}
  \caption{\textbf{Longitudinal retention.} Patients assessed at baseline and retained
  at the $12$- and $24$-month follow-up visits, overall and by cohort. Attrition is
  informative (retention depends on baseline coordinates) and is addressed throughout by
  inverse-probability weighting.}
  \label{fig:consort}
\end{figure}

\begin{figure}[hbtp]
  \centering
  \includegraphics[width=0.72\linewidth]{edfig_invariance.png}
  \caption{\textbf{Longitudinal measurement invariance.} Tucker congruence of each
  backbone axis against baseline at $12$ and $24$ months. All four backbone axes
  (\Gfac, cognition, immunometabolic, sleep) remain invariant ($\varphi\ge0.95$); the
  merged immunometabolic biology axis is fully invariant ($\varphi=0.987$).}
  \label{fig:invariance}
\end{figure}

\begin{figure}[hbtp]
  \centering
  \includegraphics[width=0.85\linewidth]{edfig_repbench.png}
  \caption{\textbf{The eight-factor map is a sufficient representation of the $143$
  raw indicators.} Held-out AUC for a matched gradient-boosting classifier predicting
  two-year deterioration and recovery, using the reference (DSM-5 $+$ severity $+$
  baseline), the reference $+$ the eight-factor map, or the reference $+$ all $143$ raw
  indicators. The map equals the raw indicators for deterioration (sufficient) and
  trails by $\approx0.04$ AUC for recovery (near-sufficient); $\ge92\%$ of the raw recovery
  signal lies within the eight factors.}
  \label{fig:repbench}
\end{figure}

\begin{figure}[hbtp]
  \centering
  \includegraphics[width=0.95\linewidth]{edfig_loso.png}
  \caption{\textbf{Leave-one-site-out external validity: the map and its headline
  finding survive out-of-site refitting.} Each row is one of the 15 sites
  ($N\ge100$) held out entirely; the GLLVM is refit on the remaining sites and the
  held-out site's loadings are projected. \textbf{(a)}~The immunometabolic loading
  pattern is essentially identical whichever site is removed (Tucker
  $\varphi=0.999$--$1.000$, all far above the $0.95$ congruence bar). All eight
  factors cleared their pre-set congruence bar in all 15 folds; the single
  weakest factor was $\varphi=0.917$ (Monaco, $n=237$), a thin two-to-four-item
  axis still above the $0.85$ bar for thin factors. \textbf{(b)}~The score-level
  dissociation between the immunometabolic axis and general burden holds in
  every fold ($|r|=0.073$--$0.082$, versus $\approx0.08$ on the full sample)---the
  finding is not driven by any single site.}
  \label{fig:loso}
\end{figure}

\begin{figure}[hbtp]
  \centering
  \includegraphics[width=0.95\linewidth]{fig_loading_vs_info.png}
  \caption{\textbf{Loading $\neq$ information: the model weighs items by what they
  reveal, not by how strongly they load.} \textbf{(a)}~Every indicator's loading on
  its home axis (on that likelihood family's own link scale, so cross-family
  magnitudes are not directly comparable) against the exact Fisher information it
  contributes to that axis score at the population mean---the quantity that
  \emph{is} comparable across families. Continuous markers with moderate loadings
  dominate; the two substance-use flags (alcohol, cannabis) load $\approx0.92$--$0.96$
  yet contribute only $\approx0.01$--$0.013$, roughly a hundred-fold less than a
  mid-loading continuous marker such as body-mass index (specific-factor loading
  $0.908$, information $2.8$; the table-reported general loading is $\approx0.95$).
  \textbf{(b)}~The mechanism: for a binary item, Fisher information is
  $\lambda^2 p(1-p)$, which collapses toward zero as endorsement approaches the
  extremes whatever the loading---so a rare binary flag is nearly uninformative
  regardless of how strongly it loads. This is why the substance-use axis, nominally
  carried by two high-loading flags, is in fact near-empty of information, and why an
  information-based reading of the map, not a loading-based one, is the correct lens
  on a mixed-type measurement model.}
  \label{fig:loadinginfo}
\end{figure}

\begin{figure}[hbtp]
  \centering
  \includegraphics[width=0.72\linewidth]{edfig_robustness.png}
  \caption{\textbf{The functional forecast is robust and course-dependent.} Held-out
  incremental value of the archetype phenotype for two-year functioning under attrition
  reweighting, cohort hold-outs and a permutation null. The signal survives attrition and
  dropping either DR or SZ, vanishes under permutation, and weakens when bipolar disorder
  is removed---it is carried by the episodic, open-course cohorts.}
  \label{fig:robustness}
\end{figure}

\clearpage
\subsection{The claims ledger}

Table~\ref{tab:claims} states, side by side, what the study claims and what it
deliberately does not---the spine of the Limitations and of any response to reviewers.

\begin{table}[hbtp]
\centering
\caption{\textbf{What we claim, and what we do not.}}
\label{tab:claims}
\renewcommand{\arraystretch}{1.3}
\begin{tabular}{p{0.46\linewidth} p{0.46\linewidth}}
\toprule
\textbf{We claim} & \textbf{We do not claim} \\
\midrule
A validated eight-dimension transdiagnostic measurement map (internal validity). &
Natural biotypes or subtypes; that no biotype exists in any measurement space. \\
Immunometabolic load \emph{largely independent} of a general functional-burden
axis (quantified, with boundary conditions). &
Strict statistical orthogonality, or independence from \emph{being ill}
(case--control elevations). \\
A graded continuum, resolved by five stable archetypes, that is descriptively
tighter than DSM-5. &
That the regions are clinically superior to DSM-5 beyond the modest prognostic
increment shown. \\
Temporal coherence: durable biology, moving symptoms. &
That developmental risk is genuinely ``state'' (it partly reflects recall noise). \\
A modest, \emph{group-level} incremental prognosis for \emph{functioning}, carried by
the biological phenotype, course-dependent. &
A large individual-risk gain; prediction of \emph{recorded} relapse/hospitalization
events; a deployable individual rule. \\
Convergent validity with the immuno-metabolic literature. &
External or out-of-sample validation (none yet). \\
\bottomrule
\end{tabular}
\end{table}
